%====================
% PROFESSIONAL EXPERIENCE
% ATS NOTE: this content MUST sit under a canonical "Experience" heading
% in sunil_resume.tex. Workday / Taleo / iCIMS / SuccessFactors segment a
% resume by standard headings to build the employment record. Without one
% they extract ZERO jobs and ZERO years, and portal filters auto-reject.
%====================
\documentclass[letter,10pt]{article}
\usepackage{TLCresume}
\usepackage{hyperref}
\begin{document}

% NOTE: do NOT wrap the title in \skills{} -- that macro is defined as
% "{ \bfseries #1}" with a LEADING SPACE, which renders as a visible indent
% inside a \subsection. The subsection titleformat already applies \bfseries.
\jobentry{Senior Data Scientist}{Aria Intelligent Solutions}{Dec 2023 -- Present}
\subsection{Agentic AI Platform -- Orchestrator \& Autonomous Data-Engineering Agents (A2A)}
\begin{zitemize}
\item Architected an \textbf{Orchestrator Agent} on the \textbf{A2A protocol} with LLM-driven \textbf{task decomposition}, knowledge-base retrieval, and \textbf{human-in-the-loop} clarification that dynamically routes to specialized sub-agents (IngestIQ, VisionIQ, InsightIQ) --- \textbf{98\%} end-to-end success across \textbf{500 concurrent sessions}, new agents onboarded with zero system-level refactoring.
\item Built \textbf{IngestIQ}, a suite of autonomous data-engineering agents (\textbf{CrewAI} + \textbf{PydanticAI} + \textbf{Claude Sonnet 4.6}) that plan and run end-to-end ETL over \textbf{databases and files} from a single natural-language task, cutting manual data-engineering effort by \textbf{90\%}.
\item Designed a \textbf{confidence-driven planner/executor} in which the LLM authors per-step \textbf{DuckDB / Python} (validator-gated, sample-first) rather than hard-coded recipes, making \textbf{DuckDB} a \textbf{cross-source streaming glue} that joins \textbf{DB tables with S3 files} and streams \textbf{Parquet} straight to S3 --- scaling to \textbf{millions of rows}, never materialized locally.
\item Hardened against \textbf{prompt-injection / jailbreak} by running all LLM-authored code in an \textbf{OS-level sandbox} (\textbf{bubblewrap}) --- scratch-confined filesystem, an \textbf{egress broker} to scoped \textbf{S3} only, \textbf{session-scoped STS credentials}, a \textbf{read-only DB gate}, and a Python \textbf{AST allow-list}.
\item Hardened for production: \textbf{human-in-the-loop} pause/resume over a \textbf{Redis} bridge, \textbf{cross-pod ownership routing} for cancel/resume, real-time \textbf{agent-brain progress streaming} via \textbf{A2A}, plus a sandboxed \textbf{content-addressed discovery cache} and any-to-any \textbf{format conversion / merge} with \textbf{bash} backends.
\item \textbf{Stack:} Python, CrewAI, PydanticAI, Claude Sonnet 4.6, A2A Protocol, DuckDB, FastAPI, Redis, Celery, PostgreSQL, AWS ECS, Athena \& S3.
\end{zitemize}

\subsection{Agent Evaluation Service -- Automated Grading \& Root-Cause Analysis}
\begin{zitemize}
\item Built an end-to-end \textbf{evaluation framework} (FastAPI \texttt{/evaluate}) that runs a registry of \textbf{28 graded tasks} --- real-world \textbf{Kaggle} datasets with ground truth from published reference notebooks --- against the \textbf{Orchestrator Agent} or any individual \textbf{A2A agent}, replacing manual QA with \textbf{regression visibility} on every release.
\item Designed a \textbf{cell-level precision / recall / F1 grader} (order-insensitive, float/date tolerant) plus an \textbf{LLM root-cause analyzer} (Claude on AWS Bedrock) adjudicating genuine errors vs.\ acceptable spec-variants; benchmarked the fleet at \textbf{86\% exact-pass} (24/28) and \textbf{0.86 mean F1}.
\item Engineered an \textbf{LLM auto-responder} driving the agent's \textbf{human-in-the-loop} flow over WebSocket for fully unattended runs, over \textbf{file- and database-backed} datasets (TPC-H, M5, Olist, Favorita) with a \textbf{Redis-snapshotted run store}.
\item \textbf{Stack:} Python, FastAPI, WebSockets, CrewAI, Claude (AWS Bedrock), Pandas, Redis, AWS S3 \& Athena.
\end{zitemize}

% Keep this heading with at least its first few bullets -- otherwise it strands
% alone at the bottom of page 1 with all its content on page 2.
\needspace{5\baselineskip}
\subsection{Content Extraction \& Tagging Pipeline}
\begin{zitemize}
\item Built an end-to-end NLP + CV pipeline parsing \textbf{millions of scientific PDFs} with \textbf{Docling} and \textbf{PyMuPDF}, robust to noisy layouts and multi-page scans.
\item Deployed \textbf{Phi-3 Vision} via \textbf{vLLM} to interpret visual content (graphs, charts), achieving \textbf{20\,ms} median inference latency in production.
\item Self-hosted \textbf{LLaMA 3 Instruct (FP8)} via \textbf{vLLM} for \textbf{hierarchical taxonomy tagging}, with a \textbf{BM25} first-pass filter (lemmatized n-gram matching of tag names, definitions and synonyms against document chunks) that prunes each taxonomy level to the top-5 candidates before the LLM is called --- far fewer LLM calls per document, proprietary data kept in-house.
\item Compressed and migrated a production \textbf{Qdrant} hybrid-retrieval index (\textbf{BGE-M3} dense, \textbf{ColBERT} late-interaction multivectors, \textbf{bm42} sparse, \textbf{CLIP} image vectors) \textbf{8x} via \textbf{TurboQuant 4-bit} at \textbf{nDCG/Recall parity} across \textbf{8 BEIR datasets} --- id-preserving migration of all 14 collections onto a separate quantized container, and root-caused \& fixed a production write-time \textbf{OOM} incident.
\item Benchmarked translation across \textbf{52 languages}; deployed \textbf{MBART-50} for multilingual inference via \textbf{FastAPI}.
\item \textbf{Stack:} PyTorch, vLLM, HuggingFace, FastAPI, PyMuPDF, Qdrant, AWS, LangChain.
\end{zitemize}

\jobentry{Applied Scientist I}{Amazon Development Centre (India) Pvt. Ltd.}{Oct 2022 -- Jun 2023}
\subsection{Ad Moderation -- Provocative Content Detection}
\begin{zitemize}
\item Developed a \textbf{FLAVA}-based multimodal classifier for inappropriate ad content detection, reaching \textbf{72\% precision} and \textbf{90\% recall}; introduced \textbf{active learning} loops targeting hard samples to lift edge-case recall and reduce manual review escalations.
\item Performed feasibility analysis and curated large-scale annotated datasets to improve model robustness across ad categories.
\item Automated cross-taxonomy mapping between IAS and internal moderation labels using Pandas pipelines over 1 year of ad data, eliminating manual reconciliation and accelerating moderation throughput through cross-functional collaboration.
\item \textbf{Stack:} PyTorch, OpenCV, HuggingFace, Pandas.
\end{zitemize}

\jobentry{Computer Vision Engineer}{Beltech AI Pvt. Ltd.}{Jul 2021 -- Sep 2022}
\subsection{Smart Attendance System -- Face Recognition}
\begin{zitemize}
\item Built and deployed a face-recognition system using fine-tuned \textbf{ArcFace} on a curated domain-specific dataset, achieving \textbf{99\% precision} and \textbf{95\% recall} on a \textbf{1{,}000-person} gallery in production.
\item Improved accuracy with keypoint-based filtering and \textbf{DBSCAN} clustering for identity resolution.
\end{zitemize}

\subsection{Multi-Camera Object Tracking \& Re-Identification}
\begin{zitemize}
\item Replaced \textbf{DeepSORT}'s default feature extractor with \textbf{ResNet-34} trained under \textbf{ArcFace metric-learning loss}, improving cross-camera re-identification across \textbf{10--20 test scenarios}.
\item Optimized assignment via the \textbf{Hungarian algorithm} and tuned \textbf{Kalman-filter} bounding-box predictions to reduce trajectory fragmentation and ID switches.
\item \textbf{Stack:} PyTorch, OpenCV, FastAPI.
\end{zitemize}

\end{document}
